\begin{table}[t]
  \centering
  \caption{%
    Calibration comparison: deterministic (STLCG) vs.\ probabilistic (pdSTL)
    planning on the single-shot scenario ($N=1000$ MC trials,
    true process noise $\sigma_q = 0.03$\,m).
    Both methods produce nominally valid plans ($\eta > 0$), confirming the
    difference is not in nominal path quality but in noise handling.
    STLCG optimises $\eta$ and provides no probabilistic claim
    (\emph{n/a} for $\rho_{\downarrow}$ at planning time); its
    $\rho_{\downarrow}$ evaluated post-hoc under true noise is low.
    pdSTL directly optimises $\rho_{\downarrow} \leq P(\varphi)$,
    and the planning-time bound tracks the deployed $\rho_{\downarrow}$
    and empirical MC success rate, demonstrating a \emph{calibrated} guarantee.
  }
  \label{tab:single_shot_comparison}
  \begin{tabular}{lcc}
    \toprule
    \textbf{Metric} & \textbf{STLCG (det.)} & \textbf{pdSTL (ours)} \\
    \midrule
    $\eta$ (nominal signed-distance) & $0.200$ (optimized) & $0.199$ (post-hoc) \\
    $\rho_{\downarrow}$ (planning-time guarantee) & n/a & $0.930$ \\
    $\rho_{\downarrow}$ deployed ($\sigma_q{=}0.03$\,m) & $0.834$ & $0.930$ \\
    MC success rate ($N=1000$) & 87.2\% & 97.8\% \\
    MC safety rate & 87.7\% & 98.4\% \\
    Min.~obstacle clearance (m) & $0.153 \pm 0.122$ & $0.259 \pm 0.097$ \\
    \bottomrule
  \end{tabular}
\end{table}
